% Nouveau EGA V (cote 25): archive cote 25, site page 172 (PDF page 173).
% typescript. Draft compared: none.
% AI-transcribed by Claude Opus 5.5 (claude-opus-5-5, adaptive thinking, effort high), Message Batches API, batch ega25b (msgbatch_01A7n5YuXubUPNK1zsepi8RG), 2026-11-13; not yet checked by eye.
\scanpage{25/172}
\mspage{7}

EGA IV 19

Remords à ``Ensemble singulier et supersingulier ...'' .

Venant de rédiger le\xout{s} sorite de la transversalité, je m'aperçois que l\unsure{e}
contexte adopté dans mes notes sur les \unsure{ensembhes|ensembles} singuliers et \unsure{supersinguliers}
savoir les sections des Modules loc libres i.e. des bons fibrés vectoriels,
n'est pas encore le ``bon''. Si $u$ est une section d'un fibré vectoriel
$E = W(\underline{E})$ (associé à $\underline{E}$ loc libre de type fini) sur X lisse sur Y, alors $V(u)$
n'est autre que l'image inverse \unsure{du} \deleted{la} \inserted{sous-préschéma $E_{o}$ de $E$} ``section nulle'', et un point $x\in V(u)$
est non singulier \inserted{pour u} rel. à S sss u est \uline{transversal} à $E_{o}$ en x \unsure{rel.} à Y.
Si cette dondition n'est pas vérifiée, i.e. si x est singulier pour u rel à
Y, nous étudions les conditions moyennant lesquelles x n'est pas trop singu-%
lier rel.à Y, i.e. u ne dévie pas trop de la transversalité en x. Je n'ai
aucun doute que tous les développements de mes notes, avec des changements
essentiellement triviaux, s'étendent au cas où on a un X-préschéma E, un
sous-préschéma $E_{o}$ de E, une section \xout{\unsure{dx}} u de E sur X, et qu'on considère l'ima-%
ge inverse $X_{o}$ de $E_{o}$ par u, \inserted[typed]{\unsure{On} suppos\xout{ant} E et $E_{o}$ lisses sur X en les points de $v(X)$.} il est possible (mais \unsure{seul} une rédaction détail-%
lée permettra de le déceler) qu'il faille supposer parfois que la dimension
relative de $E_{o}$ sur X en les points de $v(X_{o})$ soit nulle, i.e. que $E_{o}$ est
même étale sur X en ces points (ce qui sera le cas lorsque $E_{o}$ provient comme
ci-dessus d'une section de E sur X, soit $v_{o}$). Ainsi, x sera un point
``\xout{\illegible{1}} \inserted{de non transversalité à $E_{o}$} ordinaire pour v, relativement à Y'' si et seulement si \deleted{\illegible{1}} \inserted{x} est
géométriquement un point quadratique ordinaire de $X_{o\,y}$ , et si $\mathrm{codim}_{x}(X_{oy},X_{y})$
$= c$ , où c est la dodimension de $E_{o}$ dans E (en $v(x)$) ; cela implique que
$X_{o}\longrightarrow X$ est \inserted{\unsure{un immers.|une immers.}} transversalement régulier\xout{e} en x rel. à Y .... Le rôle du Module
$\underline{E}\otimes\underline{O}_{V(u)}$ sera ici joué par $v_{o}^{*}(\underline{N}_{E_{o}/E})$ , où $v_{o}:X_{o}\longrightarrow E_{o}$ est induit par v. Je
pense que tu n'auras pas de mal à tout écrire dans ce contexte,bien plus
naturel somme toute - quitte à faire à la fin les traductions dans le cas
des fibrés vectoriels. S'il y a des difficultés de démarrage déjà, nous
pourrions les travailler ensemble à ton prochain passage à Paris ... \hand{Je vais t'envoyer des notes succintes ces jours-ci.}
\margin{172}
